\section{Discussion}
\label{sec:discussion}
% =====================================================================

\paragraph{Ce qui est robuste.} La comparaison a été faite avec trois estimateurs de dérivées (émulateur, régression sur les paires, différences finies de Quijote), deux covariances, quatre HOD fiducielles, deux populations de traceurs (galaxies de type LRG, halos) et plusieurs vecteurs de données. Deux conclusions n'ont jamais changé : l'ASTRA split et le density split portent une information comparable (rapport d'erreurs 0.75--1.2 pour le même vecteur de données, avec quatre ou six paramètres, dans l'analyse de type DESI avec la covariance indépendante, 0.6--2 dans les prévisions fiables de Quijote), et ils sont complémentaires, leur combinaison gagnant de 23\% à un facteur 2.4. Le classement pour les paramètres individuels a changé avec les choix d'analyse (Tableau~\ref{tab:evolution}), au niveau de 10--40\%, ce qui est la taille de l'effet de ces choix ; nous ne proclamons donc pas de gagnant. Là où une tendance est constante, le density split est légèrement en avance sur $\omega_b$, $\Omega_m$ et $n_s$ avec les corrélations croisées des quartiles aléatoires, et l'ASTRA split en avance sur $\sigma_8$ grâce à ses auto-corrélations d'environnement et aux environnements des traceurs (Quijote). En particulier, l'avantage de l'ASTRA split sur $\sigma_8$ trouvé avec la covariance à boîte unique disparaît avec la covariance indépendante.

\paragraph{Ce que nous avons appris sur la méthodologie.} (i) Les suites d'émulation dans lesquelles la HOD est tirée indépendamment dans chaque cosmologie ne permettent pas de différences finies naïves : la discordance de HOD domine la réponse cosmologique d'un ordre de grandeur, et les différences appariées n'annulent pas le bruit d'échantillonnage des galaxies de catalogues de HOD différentes. (ii) Les covariances à boîte unique, même avec de nombreux sous-volumes et lignes de visée, peuvent biaiser la \emph{comparaison} des statistiques, et pas seulement leurs erreurs absolues. (iii) Les statistiques d'environnement de traceurs épars ont des dérivées bruitées, qu'une matrice de Fisher à six paramètres amplifie le long de ses directions dégénérées ; comparer deux estimateurs de Fisher qui encadrent la vérité permet d'identifier les prévisions fiables. (iv) Les statistiques de clustering répondent essentiellement à une seule combinaison des 12 paramètres de HOD.

\paragraph{Limitations.} (i) Les prévisions concernent un volume de $(500\Mpch)^3$, où les erreurs sont comparables aux largeurs des priors : elles comparent le contenu en information et ne sont pas des contraintes réalistes. (ii) L'inférence est calibrée mais pas encore discriminante : à ce volume le prior domine, et pour une boîte complète c'est l'erreur de l'émulateur. (iii) Les halos de Quijote sont dix fois plus épars que les LRG, et l'échelle adaptative d'ASTRA y est proportionnellement plus grande. (iv) La suite d'émulation n'a qu'une seule réalisation de HOD par point, dont le bruit d'échantillonnage des galaxies fixe le plancher de l'erreur de l'émulateur.

\paragraph{Améliorations.} Une courbe d'apprentissage montre que l'erreur de l'émulateur a atteint un plateau à environ la moitié de l'erreur statistique d'une boîte de $2\,h^{-1}$Gpc : davantage de catalogues d'entraînement du même type n'aideraient guère. Plusieurs réalisations de HOD (graines) par point, ou des catalogues appariés en HOD générés avec AbacusHOD dans toutes les cosmologies, sont nécessaires à la fois pour abaisser ce plancher, ce qui rendrait discriminante une inférence sur un volume de type DESI, et pour obtenir des dérivées propres dans le cadre de type DESI. De plus grands ensembles de boîtes indépendantes supprimeraient la dépendance restante à la covariance.

\paragraph{Au-delà des quartiles.} Locale et exempte de lissage, ASTRA fournit plus qu'un classement en densité : une classification de chaque galaxie en vides, feuillets, filaments et nœuds, une valeur d'environnement continue utilisable comme marque, et un traitement naturel des masques du relevé grâce aux points aléatoires. Le premier test des classes sur Quijote (Section~\ref{sec:classes}) est le résultat le plus prometteur de ce travail pour ASTRA : malgré leur déséquilibre, les classes avec leurs auto-corrélations contraignent $\sigma_8$ deux à trois fois mieux que le density split, de façon robuste selon les compressions. Les prochaines étapes naturelles sont les classes probabilistes de \cite{ASTRA} (classification répétée sur de nombreux catalogues aléatoires, ce qui donne à chaque point une probabilité d'appartenir à chaque classe et pourrait rendre utilisables les classes de vides et de nœuds), des seuils adaptés à la densité de l'échantillon, et un émulateur des statistiques de classes pour DESI.

